\subsection{Proper-homology detection and the first-cover normalizers}

\begin{statement}{A.11}{Proved}{A noncuspidal pro-prime normalizer is visible in proper homology}
Let $Y$ be an affine hyperbolic curve over a number field, with smooth
compactification $\overline Y$. Let $P=\pi_1(Y_{\overline K})^{(\ell)}$.
If a nontrivial closed procyclic subgroup $A\subset P$ has arithmetic
normalizer with open image and is not contained in cusp inertia, then
\[
                  A\longrightarrow V_\ell\operatorname{Jac}(\overline Y)
\]
has nonzero image. Its rational span is a line stable under an open
arithmetic subgroup. Consequently, if finitely many morphisms
$Y\to E_j\setminus\{O_j\}$ detect every potentially stable line in
proper homology, and each target satisfies the pro-$\ell$ cuspidal
normalizer criterion, then every arithmetically normalized procyclic
subgroup of $P$ is cuspidal.
\end{statement}
\begin{proof}
Choose free pro-$\ell$ generators consisting of $2g$ handle generators
and $r-1$ cusp generators, with weights one and two respectively.
The associated weighted Magnus filtration is separated: it is the
restriction of weighted degree in the noncommutative power-series ring,
and every generator has positive weight. Its rational graded pieces
are generated by the proper first homology in weight one and the
boundary Tate classes in weight two. At a good prime with residue
cardinality $q$, every eigenvalue in degree $d$ therefore has complex
absolute value $q^{d/2}$. This remains true after finite extension and
after splitting the cusp set. The construction is Galois invariant:
its degree-two subgroup is the kernel of the map to proper homology,
and commutators add weights.

Let $u$ generate $A\simeq\Z_\ell$ and let $d$ be its first nonzero
weighted degree. The leading Magnus term is nonzero in a torsion-free
module. If $\theta$ is the conjugation character of the arithmetic
normalizer, geometric inner automorphisms act trivially on that leading
term; cancellation shows that $\theta$ factors through the arithmetic
image. The resulting line in degree $d$ has character $\theta$.

By \citet[Lemma 2.1.4]{Nakamura1994}, a noncuspidal element is detected
outside the boundary in the mod-$\ell$ homology of a finite cover.
If necessary replace $u$ by a nonzero $\ell$-power and replace $P$ by
an open subgroup, as in the one-cusp case of that lemma. A nonzero
power of a noncuspidal element is still noncuspidal: otherwise its
centralizer would put the original element in the same inertia group.
The detector is an open subgroup of a pro-$\ell$ group, so its
finite cover has an $\ell$-group Galois closure and its pro-$\ell$
topology is induced. Its proper integral homology is torsion-free, so its
nonzero mod-$\ell$ image gives a nonzero rational image. Descend that
finite cover to a finite number field and intersect with the original
normalizer. Its arithmetic image is still open. The same character
$\theta$ now occurs in proper homology and has Frobenius weight one.
At one prime good for both curves, comparison gives $d=1$. By the
definition of the weighted filtration, the original image in proper
homology is nonzero, as asserted.

For the last assertion, a noncuspidal $A$ would give such a line.
One of the displayed morphisms has nonzero proper-homology image on
it, so the image of $A$ in the target pro-$\ell$ group is nontrivial.
The image of its arithmetic normalizer still has open arithmetic
projection. The target criterion puts this image in cusp inertia,
whose proper homology is zero, a contradiction. These are actual
morphisms of punctured curves; a homomorphism of Jacobians alone
would not justify this argument.
\end{proof}

\begin{statement}{A.12}{Proved}{The two CM elliptic targets}
For an elliptic curve with geometric full CM
by $\Z[i]$ or $\Z[\sqrt{-2}]$, every nontrivial closed procyclic
subgroup of the geometric pro-$\ell$ group of its once-punctured curve
with open arithmetic normalizer is contained in cusp inertia.
This holds for every prime $\ell$. In the second CM case the split
prime $3$ uses the additional rational nonvanishing proof A.21 below.
\end{statement}
\begin{proof}
All assertions are unchanged by finite extension of the field of
definition. We may work with the CM eigenspaces after an extension
defining all endomorphisms.
At a nonsplit prime, including the ramified prime $2$, proper Tate
homology has no rational line stable under an open Galois subgroup.
Indeed an ordinary good Frobenius generates the quadratic CM algebra
over $\Q_\ell$, and so does every positive power: its two eigenvalues
have quotient of infinite order, since their ideals have different
support at the two primes over its residue characteristic. A stable
line would give a unital map from that quadratic field to $\Q_\ell$.
Statement A.11 now proves the criterion at these primes.

At a split prime a stable proper line must be one of the two CM
eigenlines. We recall the rational obstruction to either line, giving
the character separation explicitly because CM by $\Z[\sqrt{-2}]$
has no automorphism of order four. Put $L=\mathbb L_{\Q_\ell}(x,y)$,
$\omega=[x,y]$, and let $b_w$ denote Lyndon bracketing with $x<y$.
In the special derivation algebra set
\[
\begin{aligned}
 \epsilon_4(x)&=\ad_x^4(y),\\
 \epsilon_4(y)&=[y,\ad_x^3(y)]-[[x,y],\ad_x^2(y)],\\
 T&=[F,\epsilon_4],\qquad F(x)=y,\quad F(y)=0.
\end{aligned}
\]
The direct tensor calculation in \texttt{a\_normalizer\_certificates.py}
gives
\[
 T(\omega)=0,\quad
 T(x)=-2b_{xxxyy}-b_{xxyxy},\quad
 T(y)=2b_{xxyyy}+3b_{xyxyy}.                    \tag{A.15}
\]
In particular $T(x)\notin[x,L_4]$ and $T(y)\notin[y,L_4]$.
For example the relevant bidegree of $[x,L_4]$ is the one-dimensional
span of $b_{xxxyy}$, while the coefficient of $b_{xxyxy}$ in $T(x)$
is $-1$. This is a rational obstruction, so it also applies to a
nonprimitive integral eigenvector.

The degree-four special outer module has the three distinct CM
characters $\psi^3\bar\psi$, $\psi^2\bar\psi^2$, and
$\psi\bar\psi^3$; its middle line is $\Q_\ell T$.
The special modules in degrees one and three vanish, and the
degree-two special module is inner, generated by $\ad\omega$.
These assertions follow either by the boundary equation
$[D(x),y]+[x,D(y)]=0$ or by its Lyndon calculation.
Nakamura's metabelian formula
\cite[Theorem A, Corollary 4.12]{Nakamura1995} identifies the middle
degree-four coordinate with a nonzero rational scalar times the
elliptic Soul\'e character $\kappa_{(2,2)}$. For split $\ell\ge5$
this character is nonzero by \cite[Theorem 4.5 and Corollary 4.10]{Ishii2025};
the fields have class number one, the multiplication is by the full
integer ring, and split primes are ordinary after finite extension.
For $\Z[\sqrt{-2}]$ at $3$, A.21 supplies the same rational assertion.
For $\Z[i]$ there is no split prime below five.

Here is why the nonzero middle coordinate gives the needed individual
derivation even if the two outside coordinates are also present.
The rational span of the degree-four Galois image is stable under the
CM torus. The ratio $\psi/\bar\psi$ has infinite image on every open
Galois subgroup, so the three displayed characters remain distinct.
Choose an element separating their three eigenvalues; polynomial
projectors in its action show that a nonzero middle projection puts
the entire line $\Q_\ell T$ in this span. One can compute the image
on the kernel of the homological representation. All lower outer
degrees are zero; their possible degree-two inner term can be
removed by a rational power of the actual boundary element, whose
logarithm starts with $\omega$. This keeps the boundary fixed, so
the special degree-three vanishing still applies. Thus this is a statement in the
actual first non-inner outer quotient, not merely in a universal
space of special derivations. Nonvanishing persists on every open
subgroup, since an additive $\ell$-adic character with nonzero image
cannot have finite image.

Suppose a procyclic subgroup with first term $x$ were normalized
arithmetically. An element in the kernel of the homological action
must act on this subgroup with exponent one. After the inner
normalization just described, its first non-inner action is in
degree four. Writing a proposed generator as $x+$ higher terms in
the rational nilpotent Lie algebra through degree five, invariance
up to conjugation forces $D(x)\in[x,L_4]$. To justify the degree of
the conjugator, cancel its successive components proportional to
the proposed generator; the free-Lie centralizer of $x$ is
$\Q_\ell x$, so its first effective component has degree four.
Contributions of the
higher terms start in degree six. This is a linear condition on
$D$, so it holds on the rational span of the Galois image. It
contradicts (A.15) for $D=T$. The other eigenline is excluded in
the same way. A.11 now finishes the proof in the original pro-$\ell$
group. We have used the Malcev quotient only for a necessary
condition, not for a claim about integral conjugacy.
\end{proof}

\begin{statement}{A.13}{Proved}{Punctured elliptic maps for the quaternion and stable dihedral covers}
Work over an algebraic closure of $\mathbf Q(i)$; every curve, map,
origin, and isogeny below is defined over a finite extension.
Let $S$ be the cusp set specified in A.4.

For the quaternion cover
\[
 C_Q:X^4+Y^4=2Z^4,\qquad S=\{Z=0\},
\]
each of the three sign-involution elliptic quotients can be followed
by an isogeny so that the inverse image of its target origin is
\emph{exactly} $S$. For the sign changes of $X$ and $Y$ these
composites have degree four and have target
$E':B^2=A^3-32A$, with affine formulas
\[
 \phi_X(x,y)=(4y^2,-8i\,yx^2),\qquad
 \phi_Y(x,y)=(4x^2,-8i\,xy^2),\qquad x=X/Z,\quad y=Y/Z.
                                                        \tag{A.13.1}
\]
Their fibres over $O$ are reduced. For the sign change of $Z$,
choose one cusp image as origin on the elliptic quotient and
follow its degree-two quotient map by multiplication by two.
The resulting degree-eight map has fibre $2\sum_{P\in S}P$
over the origin. All three elliptic targets have CM by
$\mathbf Z[i]$, and their pulled-back differential lines span
$H^0(C_Q,\Omega^1)$.

For the stable dihedral genus-two curve
\[
 D_c:z^2=(t^2-4c)(t^4-4ct^2+2c^2),\qquad c^4=-1,
\]
write $Q_+,Q_-$ for its two points at infinity. There is a
degree-two map
\[
 D_c\longrightarrow E_c^{-2}:\quad (t,z)\longmapsto(X,Y)=(t^2,z),
 \qquad
 Y^2=X^3-8cX^2+18c^2X-8c^3.                         \tag{A.13.2}
\]
This elliptic curve has $j=8000$ and CM by
$\mathbf Z[\sqrt{-2}]$. The origin fibre is precisely
$Q_++Q_-$. It is the complementary elliptic factor to the
$\mathbf Q(i)$ factor of A.4 and defines a punctured map with
no extra points removed from the source.

For the stable symmetric-group genus-two curve
\[
 D_3:z^2=(t^2+4)(t^3+3t),\qquad P=\infty,
\]
every nonconstant morphism to an elliptic curve whose origin fibre
is supported only at $P$ has degree three and ramification divisor
$2P$. All such maps pull back elliptic differentials to the same
one-dimensional subspace of $H^0(D_3,\Omega^1)$; they cannot detect
the second elliptic factor independently. In fact they coincide
up to an isomorphism of their pointed elliptic targets.
Thus the proposed second independent punctured elliptic map
for this one-cusp genus-two curve does not exist.
\end{statement}
\begin{proof}
For the sign change of $X$, the quotient has coordinates
$t=Y/Z$, $u=X^2/Z^2$ and equation
\[
                         u^2=2-t^4.
\]
This is a smooth genus-one quartic. Its two points at infinity,
$P_+$ and $P_-$, have $u/t^2=+i$ and $-i$, respectively.
The sign involution has no fixed cusp, so the four cusps map
in pairs to these two points and are unramified there.
The identity
\[
                   (u-it^2)(u+it^2)=2
\]
and the expansions at infinity give
\[
             \operatorname{div}(u-it^2)=2P_+-2P_-.
\]
The difference $P_--P_+$ is nonzero: otherwise a function with
one simple pole would give a degree-one map from this elliptic
curve to a projective line. Choose $P_+$ as origin. The other
point is therefore an exact two-torsion point, and the isogeny
with kernel these two points has precisely their union as its
origin fibre. Pulling back gives exactly the four original cusps.

There are explicit coordinates for this isogeny. Put
$v=-iu$, so $v^2=t^4-2$, and set
\[
       a=2(v+t^2),\qquad b=4t(v+t^2).
\]
Then $b^2=a^3+8a$, the chosen origin is infinity, and the other
infinity is $(0,0)$. The standard degree-two isogeny with kernel
$(0,0)$ is
\[
              (a,b)\longmapsto
                   \left(a+\frac8a,b\left(1-\frac8{a^2}\right)\right).
\]
Its two coordinates simplify to $4t^2$ and $8tv$.
This proves the first formula in (A.13.1); interchange $X,Y$ to
obtain the second. The target has the automorphism
$(A,B)\mapsto(-A,iB)$, and is isomorphic to $y^2=x^3-x$
after adjoining a fourth root of $32$.

For the sign change of $Z$, use $t=X/Y$, $u=Z^2/Y^2$ on
the quotient. Its equation is $2u^2=t^4+1$. The four cusp
images are exactly the four points $u=0$. Choose one of these
as origin. The involution $u\mapsto-u$ fixes this origin and
is the elliptic involution $[-1]$; hence those four points are
exactly the full two-torsion subgroup. Multiplication by two
has precisely that subgroup as its origin fibre. Each original
cusp is ramified with index two in the sign-$Z$ quotient, so the
composite fibre has multiplicity two at each of the four cusps.
The quartic has $j=1728$, either by its order-four automorphism
or by the elementary quartic invariant formula. Its target under
$[2]$ is the same elliptic curve.

Finally the canonical forms on the plane quartic are the
adjunction forms corresponding to $X,Y,Z$. Under the three
sign involutions their respective invariant lines are distinct
and together span the three-dimensional differential space.
These are the pullbacks from the three genus-one quotients.
An isogeny in characteristic zero has nonzero differential,
so composing with the displayed isogenies leaves those three
lines unchanged. This proves the claimed simultaneous detection
of all three proper elliptic factors.

For (A.13.2), direct substitution $X=t^2$, $Y=z$ gives the cubic.
Its discriminant and invariant are
\[
                       \Delta=512c^6,\qquad c_4=160c^2,
\]
so it is smooth and $j=c_4^3/\Delta=8000$. At each infinity
of $D_c$, the functions $X$ and $Y$ have pole orders two and
three, respectively. Thus each point maps unramified to the
cubic's origin, and their total degree two exhausts the fibre.

The CM assertion can also be proved without consulting a table
of singular moduli. Translate $x=X-4c$ to write the cubic as
\[
                         y^2=x^3+4cx^2+2c^2x.
\]
The two-isogeny with kernel $(0,0)$ has target
$y'^2=x'^3-8cx'^2+8c^2x'$. This target is isomorphic to the
source by $(x',y')=(-2x,2\sqrt{-2}\,y)$.
Compose the two-isogeny with the inverse isomorphism. The resulting
degree-two endomorphism has differential $-\sqrt{-2}$, so its
square is multiplication by $-2$: endomorphisms of an elliptic
curve in characteristic zero are determined by their differentials.
This embeds the maximal quadratic order $\mathbf Z[\sqrt{-2}]$
in its endomorphism ring. It is therefore the CM ring.
It differs from $\mathbf Q(i)$, so this elliptic quotient is
the complementary factor in A.4.

For the last assertion, the point at infinity on the odd-degree
hyperelliptic model $D_3$ is a Weierstrass point. If
$f:D_3\to E_0$ is nonconstant and $f^{-1}(O)$ is supported
only there, then $f^*O=dP$, where $d=\deg f$. The ramification
divisor $R_f$ has degree two by Riemann--Hurwitz and contains
$(d-1)P$. Thus $2\le d\le3$. The case $d=2$ is impossible:
$R_f$ is the divisor of the pullback of a nonzero elliptic
differential, hence is canonical. The unique canonical divisor
containing the Weierstrass point $P$ is $2P$, whereas a
degree-two map has ramification multiplicity at most one at
each point. Therefore $d=3$ and $R_f=2P$.
Consequently
\[
                 f^*H^0(E_0,\Omega^1)=H^0(D_3,\Omega^1(-2P)),
\]
a one-dimensional space. The degree-three map already supplied
by A.4 gives precisely this line.

For completeness, equality of these lines gives uniqueness of
the pointed target quotient. Embed the field of definition in
$\mathbf C$. The induced homomorphism
$J(D_3)\to E_0$ has connected kernel. Otherwise $f$ factors
through a nontrivial finite étale isogeny of elliptic curves;
its degree divides three, and the remaining map would have
degree one, which a genus-two curve cannot have. The connected
kernel is the elliptic subvariety whose tangent space is the
annihilator of the displayed differential line. A connected
complex subtorus is determined by its tangent space. Hence all
these kernels are equal, and the quotient elliptic curves are
isomorphic. Since $P$ maps to the chosen origins, no translation
ambiguity remains.

All displayed functions extend to morphisms of smooth proper
curves, and the exact fibre computations show that they restrict
to morphisms of the indicated punctured curves. After compatible
geometric base points and paths are chosen, they induce the
usual continuous homomorphisms of arithmetic and geometric
fundamental groups. A change of path conjugates these maps by
a geometric element; the proper-homology detection assertions
are unaffected. The executable certificate
\texttt{code/a\_cover\_maps.py} verifies the map identities,
the two-isogeny identities, and the elliptic invariants.
\end{proof}

\IfFileExists{sections/A_prime_three.tex}{\begin{statement}{A.21}{Proved}{The diagonal elliptic Soul\'e character at the split prime three}
Let \(k=\mathbb Q(\sqrt{-2})\), let
\(3=\mathfrak p\overline{\mathfrak p}\), and let \(E/k\)
be a CM elliptic curve with endomorphism ring
\(\mathbb Z[\sqrt{-2}]\), good reduction at three, and
its standard puncture at the origin. In the homological
normalization \(\psi\overline\psi=\chi_3\), the diagonal
elliptic Soul\'e character
\(\kappa_{(2,2)}\) is nonzero after tensoring with
\(\mathbb Q_3\). Its restriction to any finite-index
subgroup of its defining Galois group remains nonzero.
This assertion at three is proved below; it is not an
application of a theorem whose statement assumes
\(p\ge5\).
\end{statement}
\begin{proof}
Write \(S\) for the two primes above three and the
infinite places. We first prove the precise cohomological
input, including the residual character:
\[
 H^2(G_{k,S},\mathbb Z_3(2))=0.                         \tag{A.21.1}
\]
Put \(L=k(\zeta_3)=\mathbb Q(\sqrt{-2},\sqrt{-3})\).
The discriminants of its three quadratic subfields are
\(-8,-3,24\); hence \(\operatorname{disc}(L)=576\).
Its Minkowski bound is \(36/\pi^2<4\). There is no
prime ideal of norm two: at two the extension supplied
by \(\sqrt{-3}\) is the unramified quadratic extension,
so the residue degree is two. The only prime ideals of
norm below four are the two primes above three. Each
has square principal, generated by
\(1+\sqrt{-2}\) or \(1-\sqrt{-2}\). Thus the class
group is generated by elements of order dividing two,
and \(\operatorname{Cl}(L)[3]=0\).

There are exactly two primes in \(S_L\), both fixed by
\(A=\operatorname{Gal}(L/k)\simeq C_2\). Kummer theory
and the Brauer invariant-sum sequence give
\[
 H^2(G_{L,S_L},\mu_3)
  =\ker(\mathbb F_3^2\xrightarrow{\sum}\mathbb F_3)
  \simeq\mathbb F_3,
\]
with trivial \(A\)-action. Passing from \(\mu_3\) to
\(\mathbb F_3(2)\) tensors this action by the nontrivial
cyclotomic character of \(A\). Averaging over \(A\)
is exact, and therefore
\(H^2(G_{k,S},\mathbb F_3(2))=0\). The coefficient
sequences prove vanishing for every
\(\mathbb Z/3^n(2)\). Finite-ramification cohomology
with finite coefficients is finite, so the
Mittag--Leffler property and continuous comparison
prove (A.21.1). This uses the same precise finiteness
and comparison theorems as D.1.

We next explain why the cyclotomic elliptic-unit map
has nonzero rational image, without any integral
claim that twelve is a three-adic unit. Put
\(L_\infty=k(\mu_{3^\infty})\),
\(\Delta=\operatorname{Gal}(k(\mu_3)/k)\simeq C_2\),
and \(\Gamma=\operatorname{Gal}(L_\infty/k(\mu_3))\simeq\mathbb Z_3\).
Let \(\mathcal E_L\) be the inverse limit under norms
of the three-adic completions of global units in this
tower, and let \(\mathcal A_L\) be the inverse limit
of the three-primary ideal class groups. A superscript
\((1)\) in the next formulas means the cyclotomic
Tate twist. The norm-and-Kummer construction gives
\[
 e_L:\mathcal E_L(1)_{\Delta\times\Gamma}
              \longrightarrow H^1(G_{k,S},\mathbb Z_3(2)).
                                                               \tag{A.21.2}
\]
Both its kernel and the twisted class-group coinvariants
are finite. Here is the cohomological verification of
the rational assertion needed, so that no restriction
\(p\ne3\) from a regulator theorem is imported.

Form the inverse corestriction complex
\(C_{\rm Iw}=R\varprojlim_nR\Gamma(G_{L_n,S},\mathbb Z_3(1))\).
Shapiro's cochain isomorphism, followed by the two-term
resolution for \(\Gamma\), identifies its derived
specialization at the character \(\chi_3^{-1}\), with
the corresponding \(\Delta\)-component, with
\(R\Gamma(G_{k,S},\mathbb Z_3(2))\).
The finite-level Kummer sequences identify its first
cohomology with norm-compatible \(S\)-units and its
second cohomology with an extension of
\(\mathcal A_L\) by the local Brauer-invariant modules.
The valuation and Brauer modules at the two primes
above three have trivial \(\Delta\)-action, whereas
\(\chi_3^{-1}|_\Delta\) is nontrivial. They therefore
vanish in this component; ordinary units and
\(S\)-units give the same component. The resulting
one-variable specialization sequences are
\[
\begin{split}
0&\longrightarrow\mathcal E_L(1)_{\Delta\times\Gamma}
 \longrightarrow H^1(G_{k,S},\mathbb Z_3(2))
 \longrightarrow\mathcal A_L(1)^{\Delta\times\Gamma}
 \longrightarrow0,\\
H^2(G_{k,S},\mathbb Z_3(2))
 &\simeq\mathcal A_L(1)_{\Delta\times\Gamma},
\end{split}                                                   \tag{A.21.3}
\]
up to finite groups, which is sufficient here. More
explicitly, any finite error from passing between
finite-level Kummer sequences and completed units
vanishes after tensoring with \(\mathbb Q_3\); the
completed cohomology groups are compact; inverse limits
of their exact sequences are exact by compactness. At
each finite coefficient level the groups are finite,
so the corresponding diagonal towers are Mittag--Leffler. The maps in (A.21.3) are the
norm-and-Kummer maps defining (A.21.2), not different
abstract comparison maps. The class-group Iwasawa
module is finitely generated torsion. For such a
one-variable module, finite twisted coinvariants imply
finite twisted invariants, by its elementary-divisor
decomposition. Equation (A.21.1) proves this finiteness
and makes (A.21.2) an isomorphism after tensoring with
\(\mathbb Q_3\). The global Euler characteristic formula
then gives
\(\dim_{\mathbb Q_3}H^1(G_{k,S},\mathbb Q_3(2))=1\).
Thus the rational image of (A.21.2) is nonzero.

Let \(\mathcal C_L^{\rm R}\subset\mathcal E_L\) be
Rubin's elliptic-unit module. Rubin's original
hypotheses hold at three: the Hilbert class field of
\(k\) is \(k\), its roots of unity have order two,
\([k(\mu_3):k]=2\), and three splits in \(k\).
In particular \cite[\S4, Theorem 4.1(i)]{DE_Rubin1991}
allows this prime; its restriction is
\(p\nmid\#\mu(H)[k_0:k]\), not \(p\nmid6\).
The first characteristic-ideal equality in that
theorem identifies
\(\operatorname{char}(\mathcal E_L/\mathcal C_L^{\rm R})\)
with \(\operatorname{char}(\mathcal A_L)\).
The finite specialization in (A.21.3), the
one-variable torsion structure theorem, and
\cite[Corollary 7.8]{DE_Rubin1991} show that
\[
 (\mathcal C_L^{\rm R}(1)_{\Delta\times\Gamma})\otimes\mathbb Q_3
 \longrightarrow
 (\mathcal E_L(1)_{\Delta\times\Gamma})\otimes\mathbb Q_3
                                                               \tag{A.21.4}
\]
is an isomorphism. The same conclusion follows by
taking the invariants and coinvariants of
\(0\to\mathcal C_L^{\rm R}\to\mathcal E_L
\to\mathcal E_L/\mathcal C_L^{\rm R}\to0\): both
end terms from the quotient are finite. Consequently
Rubin units have nonzero image under (A.21.2).

We transfer this image from the cyclotomic tower to
\(k(3^\infty)\). Write \(\mathcal C^{\rm R}\) for
its two-variable elliptic-unit module. The kernel and
cokernel of the norm map
\[
 (\mathcal C^{\rm R})_{\operatorname{Gal}(k(3^\infty)/L_\infty)}
                   \longrightarrow\mathcal C_L^{\rm R}
\]
are annihilated by the product of the annihilator of
\(\mathbb Z_3(1)\) and the augmentation ideals of the
two local decomposition groups, by
\cite[Theorem 7.7(ii)]{DE_Rubin1991}. After twist \((1)\)
and coinvariants, a displayed annihilator evaluates as
\[
 (\chi_3(\sigma)^{-1}-\chi_3(\sigma))
 (\chi_3(\tau_1)^{-1}-1)
 (\chi_3(\tau_2)^{-1}-1).                              \tag{A.21.5}
\]
The cyclotomic character maps each local decomposition
group onto an open subgroup of \(\mathbb Z_3^\times\).
Choose all three character values equal to a common
\(u=4^a\), \(a\ge1\), lying in these images. Then
(A.21.5) is \((u^{-1}-u)(u^{-1}-1)^2\ne0\).
For \(u=4\) its value is \(-135/64\).
Localization at the evaluation character therefore
annihilates the kernel and cokernel before taking
coinvariants, so no nonexact tensor step is being
silently used. The rational norm map on the required
specializations is an isomorphism. Norm compatibility
of Kummer corestriction proves that the two-variable
Rubin units also have nonzero rational Soul\'e image.

Now compare the units actually defining the geometric
character. Let \(\mathcal C^{\theta}\) be formed from
the theta values of \cite[\S2.2]{DE_Ishii2024}.
At every finite level, the definitions give
\[
 12 C^{\rm R}\ \subset\ C^{\theta}\ \subset\ C^{\rm R}
                                                               \tag{A.21.6}
\]
in additive notation, after adjoining the same roots
of unity. A Rubin generator is an appropriate twelfth
root of the corresponding theta-unit generator;
raising it to twelve gives that generator, and theta
units themselves belong to Rubin's group. These
inclusions commute with norms and use the same
integer twelve at every level. Hence (A.21.6) also
holds for the inverse-limit modules and makes their
rationalizations equal. Roots of unity cause no
additional image: after twist \((1)\) they have
character \(\chi_3^2\), whose specialization at the
trivial character is rationally zero. Thus theta
units have nonzero rational Soul\'e image at three.
This is the step for which the published argument
uses \(p\ge5\); replacing an integral equality by the
uniform inclusions (A.21.6) proves the rational
assertion at three.

Finally the distribution and Kummer identities in
\cite[\S3.1 and \S4.1]{DE_Ishii2024} identify that
image with multiples of \(\kappa_{(2,2)}\). We record
the specialization explicitly. Take
\(\pi=1+\sqrt{-2}\), so \(\pi\overline\pi=3\),
and \(w_k=2\). For any auxiliary ideal \(\mathfrak a\)
prime to six, the primitive theta-unit character
\(\kappa_{(2,2),\mathfrak a}\) satisfies
\[
 \underbrace{2\left(\frac1{1-3}+\frac1{1-3}-1\right)}_{-4}
       \kappa_{(2,2),\mathfrak a}
   =\left(N\mathfrak a-(N\mathfrak a)^{-1}\right)
       \kappa_{(2,2)}.                                  \tag{A.21.7}
\]
The identity follows by separating the two coordinates
of the primitive torsion sum, applying the theta
distribution relation, and summing the geometric
series of ratio three. Its denominators \(1-3\)
are three-adic units, so this calculation is valid at
three. The factor on the right is nonzero for every
nontrivial auxiliary ideal. The theta units in the
two-variable inverse limit are generated by the
primitive theta values and their Galois translates:
the one-coordinate units vanish in this inverse
limit, because norms along the other tower raise
those units to arbitrarily high powers of three.
Taking \(\sigma-1\) simply multiplies their specialized
Soul\'e images by \(\chi_3(\sigma)^{-1}-1\).
Thus a zero \(\kappa_{(2,2)}\) would make the entire
nonzero theta-unit image vanish, contradicting the
preceding paragraph. This proves nonvanishing.
For example \(\mathfrak a=(5)\) has norm twenty-five
and the right factor in (A.21.7) is \(624/25\ne0\).

All arguments just given concern rational nonvanishing;
no assertion about integral surjectivity is required.
The character is a continuous homomorphism to a
characteristic-zero three-adic line on its defining
kernel. If its restriction to a finite-index subgroup
were zero, its image would be finite and hence zero.
This proves the final assertion.
\end{proof}
}{}

\begin{statement}{A.14}{Proved}{The normalized procyclic subgroups in the first covers}
For every prime $\ell$, the nontrivial closed procyclic subgroups
with open arithmetic normalizer in the following pro-$\ell$ groups
are exactly the nontrivial closed subgroups of conjugates of cusp
inertia: all three $S_3$ covers $C_3$, all three dihedral covers
$C_4$, the quaternion cover $C_Q$, all dihedral genus-two quotients
$D_4$, and the two nonstable $S_3$ genus-two quotients $D_3$.
The same assertion holds for the stable $S_3$ genus-two quotient
$D_3$ at $\ell=2$. Statement A.25 below settles its odd-prime
normalizer problem by a further monodromy detector. The requested
nonabelian fixed subgroups are computed in A.22, A.24, and A.27 in the
specified cases; the remaining fixed-group problems are recorded
separately in A.10.
\end{statement}
\begin{proof}
First note that every nontrivial closed subgroup of cusp inertia
has open arithmetic normalizer, after making that cusp rational.
We prove the converse using A.11 and the explicit models of A.4.

On the stable $C_3$, its map to $E_{0,-1}$ followed by the
two-isogeny to $E$ has both cusps as its entire origin fibre.
The elliptic quotient by $hj$ also maps exactly these two cusps
to its origin. This quotient has $j=1728$. Composing it with a
cubic rotation gives a second independent elliptic differential
on the complementary two-dimensional summand. Together the three
maps detect all of proper homology. Statements A.11--A.12 apply.

For a nonstable $C_3$, the complementary elliptic factor
$u^2=t^3-3t+6$, up to twist, is non-CM. Its good ordinary
Frobenius discriminants at $5$ and $11$ are $-19$ and $-35$,
respectively, as independently counted in
\texttt{a\_normalizer\_certificates.json}.
A characteristic-zero CM field would inject into each ordinary
reduction's geometric endomorphism field, giving two different
quadratic fields, which is impossible. Similarly the nonstable
$C_4$ complement $u^2=t^4-4t^2+8$ has discriminants $-16$ and
$-12$ at $5$ and $7$, and is non-CM. Its other reflection quotient
is isogenous to it, so has the same conclusion. The sign changes
of the parameter give twists and do not change geometric CM.

An elliptic curve with geometric endomorphism ring $\Z$ has no
Tate line stable under any open arithmetic subgroup. Indeed Tate
semisimplicity and the endomorphism theorem, over that finite
extension, would give endomorphism dimension at least two if its
two-dimensional representation had an invariant line. The actual
endomorphism algebra is $\Q_\ell$. A line in a direct sum of
copies would project to such a line in one copy. Thus in each
nonstable $C_3$ or $C_4$, all potentially stable proper lines are
in the $E$ factor, which the original covering map to $U$ detects.
This proves these cases, and the same reasoning proves their
genus-two quotient cases, since the descended map to $U$ has
exactly the indicated cusps above $O$.

For stable $C_4$, the proper factors are $E$ and two copies of
the elliptic curve with $j=8000$. Its elliptic reflection quotient
is $u^2=t^4-4ct^2+2c^2$. The degree-two map
$(t,u)\mapsto(X,Y)=(t^2,tu)$ has target
$Y^2=X^3-4cX^2+2c^2X$, again with $j=8000$.
Its origin fibre is exactly the two infinity points.
Consequently the composite map from $C_4$ has the four cusps as
its entire origin fibre. Rotate the reflection to obtain the
second independent complementary differential. By A.13 these
targets have CM by $\Z[\sqrt{-2}]$, and A.12 applies at every
prime. The stable $D_4$ has the additional explicit degree-two
map of A.13 to a once-punctured elliptic curve with this same CM,
which detects its complementary factor. The original map to $E$
detects the remaining factor.

The three punctured elliptic maps of A.13 for $C_Q$ have $j=1728$
targets and span its three-dimensional differential space.
They therefore detect proper homology, proving the quaternion case.

Finally the stable $C_3\to D_3$ is an etale double covering,
also at the compactification. At $\ell=2$ its pro-$2$ group is
an open subgroup of index two of that of $D_3$; the pro-$2$
topology on an index-two subgroup is induced. A nontrivial
arithmetically normalized procyclic subgroup intersects it in
a nontrivial subgroup, which is cuspidal by the already proved
$C_3$ case. The centralizer of this nonzero inertia subgroup is
the containing inertia in $D_3$, so the original subgroup is
cuspidal as well. At odd primes this inclusion argument is not
valid. A.13 proves that no second independent elliptic map with
only the single cusp above its origin can supply the missing
detector, so the same method cannot silently be extended there.
\end{proof}
